\section{Boundary Green functions and stability on the tree}\label{sec:tree}

The local law will compare a random empirical distribution of Green
functions with the cavity Green function on the infinite tree.  We first
prove the distributional stability needed for that comparison.  The same
argument constructs boundary values in the common potential environment
and supplies the integrability needed when we later take entropies of
Gaussian fields on an edge or a star.

Weak-disorder absolute continuity on regular trees was established by
Klein \cite{Klein1998}; the stability of absolutely continuous spectral
weight was studied by Aizenman, Sims, and Warzel \cite{ASW2006}.
The geometric Green-function argument of Froese, Hasler, and Spitzer
\cite{FHS2007}, for the tree of degree three, is a precedent for using the
half-plane geometry of the recursion.  We give the quantitative estimates
used here, including their proofs, for every fixed $d\ge3$.

Fix a compact interval $J\subset(-2\sqrt b,2\sqrt b)$, where $b=d-1$.
A forward tree is the component left after removing the edge from its
root towards its parent.  Its root has $b$ children.  Write $g(\zeta)$
for its diagonal Green function, with $\zeta=e+i\eta$ and $\eta>0$.
Schur complementation gives
\begin{equation}\label{eq:tree-recursion}
 g(\zeta)=\left(V-\zeta-\sum_{i=1}^b g_i(\zeta)\right)^{-1},
\end{equation}
where the $g_i$ are the cavity Green functions of the child branches.
For independent potentials these branches are independent and have the
same law as $g$, independently of $V$.  The clean boundary solution is
\[
 q=q(e)=\frac{-e+i\sqrt{4b-e^2}}{2b},
 \qquad q=(-e-bq)^{-1},\qquad |q|=b^{-1/2}.
\]
We use the disk coordinates
\begin{equation}\label{eq:tree-disk}
 X_e(g)=\frac{g-q(e)}{g-\overline{q(e)}},\qquad
 T_e(x)=\frac{q(e)-x\overline{q(e)}}{1-x},\qquad
 D(x)=\frac1{1-|x|}.
\end{equation}
Thus $X_e:\HH\to\DD$ and $T_e=X_e^{-1}$.  When the real part is
understood we write $X=X_e$.  Put
$Y(e,\eta)=Y(e+i\eta)=X_e(g(e+i\eta))$.  For $j\ge1$, let $F_j$ be the map from
the $b^j$ disk-valued leaves to the root obtained by applying
\eqref{eq:tree-recursion} through $j$ levels and then using $X_e$.
Its other arguments are the $1+b+\cdots+b^{j-1}$ internal potentials.
Whenever $F_j$ is used as a map on probability laws, these potentials
are independent and uniform on $[-w,w]$, independently of the leaves.
The notation $\|Z\|_2$ means $(\EE|Z|^2)^{1/2}$, and $W_2$ is the
Euclidean Wasserstein distance on $\DD$.

\begin{proposition}[Tree stability]\label{prop:tree-stability}
There is an integer $\ell\ge1$, depending only on $b,J$, with the
following properties.  Given any finite order $p_*>0$ and a sufficiently
small prescribed number $\delta_{\rm tree}>0$, one can choose
$h,w_0>0$ so that, uniformly for $e\in J$, $0<w\le w_0$, and
$0<\eta\le h$,
\begin{equation}\label{eq:tree-small-moments}
 \|Y(e,\eta)\|_2\le\delta_{\rm tree},\qquad
 \EE D(Y(e,\eta))^p\le K_p\quad(0<p\le p_*).
\end{equation}
The choices also include the additional fixed finite moment orders
required by the stability assertion that follows.  There are constants
$\delta_*>0$ and $C<\infty$ such that, if a probability law $P$ on
$\DD$ satisfies $\int|x|^2\,dP(x)\le\delta_*^2$ and
\[
 W_2\bigl(P,\Law(F_j(X_1,\ldots,X_{b^j}))\bigr)\le u,
 \qquad j=1,\ell,
\]
where the $X_i$ are independent with law $P$, then
\begin{equation}\label{eq:tree-law-stability}
 W_2(P,\Law(Y(e,\eta)))\le Cu.
\end{equation}
The constants $K_p,\delta_*,C$ can be fixed before decreasing the
parameters further to make $\delta_{\rm tree}$ much smaller than
$\delta_*$.  On a fixed tree with its fixed potential variables,
$(e,\eta)\mapsto Y(e,\eta)$ extends continuously in $L^2$ to
$J\times[0,h]$.  The boundary values belong to $\DD$ almost surely,
obey \eqref{eq:tree-small-moments}, and are measurable functions of
the potentials.
\end{proposition}

The proof separates the deterministic geometry from the probabilistic
contraction.  A single level need not contract the mean, and arbitrary
disk inputs need not have uniformly bounded Euclidean derivatives.
The useful facts are that the mean multiplier is separated from $1$,
and that many levels strongly damp a perturbation arriving through one
leaf.  Independence lets us combine these two facts.

\subsection{Deterministic estimates for the recursion}

All constants in this subsection are uniform for $e\in J$ and for
sufficiently small bounded $|V|,\eta$.  At $V=\eta=0$, with all but
one child equal to $0$ in disk coordinates, the one-level map is
\[
 H(x)=\frac{ax}{b-(b-1)x},\qquad a=bq^2.
\]
Here $|a|=1$ and $\inf_{e\in J}|1-a(e)|>0$.  Along a path through
$j$ levels, all other leaves being zero, the map is $H^{\circ j}$.
Direct composition gives
\[
 H^{\circ j}(x)=\frac{(a/b)^j x}{1-c_jx},\qquad
 c_j=\left(1-\frac1b\right)\sum_{k=0}^{j-1}(a/b)^k.
\]
The limiting coefficient is $(b-1)/(b-a)$ and has modulus strictly
less than one, uniformly on $J$.  Consequently, for all sufficiently
large $j$,
\begin{equation}\label{eq:tree-one-leaf}
 \sup_{|x|\le1}|H^{\circ j}(x)|\le Cb^{-j},\qquad
 \sup_{|x|\le1}|(H^{\circ j})'(x)|\le Cb^{-j}.
\end{equation}
Fix $\ell$ large enough that the first bound is less than $1/2$ and
$2Cb^{-\ell}b^{\ell/2}<1/10$.  The maps and their derivatives remain
bounded on a neighborhood of each set with one unrestricted input
$|x_i|\le1$ and all other inputs zero.  Indeed, siblings whose disk
coordinates are close to zero have Green functions close to $q$, so the one-variable fractional-linear
map has no pole on the closed disk; the apparent singularity at
$x_i=1$ cancels after multiplication by $1-x_i$.  At successive
vertices on the path the same statement applies, while off-path
branches remain close to their clean values.  Compactness and
\eqref{eq:tree-one-leaf} therefore give, after decreasing $h,w_0$,
\begin{equation}\label{eq:tree-one-leaf-perturbed}
 \Lip\bigl(x\mapsto F_\ell(0,\ldots,x,\ldots,0)\bigr)
 \le 2Cb^{-\ell},
\end{equation}
and a strictly positive output margin whenever all but at most one
leaf have modulus at most a fixed sufficiently small $r>0$.
The derivatives needed in such neighborhoods are uniformly bounded.

Two estimates handle the remaining inputs.  For $g\in\HH$,
\begin{equation}\label{eq:tree-hyperbolic-comparison}
 D(X_e(g))\asymp\frac{1+|g|^2}{\Im g}.
\end{equation}
To verify this, use
$1-|X_e(g)|^2=4\Im q\Im g/|g-\bar q|^2$ and
$|g-\bar q|^2\asymp1+|g|^2$, uniformly in $e\in J$.
If $A=V-e-i\eta-\sum_i g_i$, then
$-\Im A=\eta+\sum_i\Im g_i$, and
\[
 \frac{1+|A^{-1}|^2}{\Im(A^{-1})}
 =\frac{1+|A|^2}{\eta+\sum_i\Im g_i}
 \le C\sum_i\frac{1+|g_i|^2}{\Im g_i}.
\]
The last inequality uses boundedness of $V-e-i\eta$ and
$|\sum g_i|^2\le b\sum|g_i|^2$.  Thus
\begin{equation}\label{eq:tree-margin-bound}
 D(F_1(y_1,\ldots,y_b))\le C\sum_iD(y_i).
\end{equation}
Its iteration bounds the output inverse margin by $C_j\sum_iD(y_i)$.
The same one-level estimate holds with $d$ children at a full root.

For each fixed $j$, there is a finite exponent $A_j$ such that
\begin{equation}\label{eq:tree-sibling-bound}
 |F_j(x)-F_j(y)|
 \le C_j\sum_{i=1}^{b^j}|x_i-y_i|
                    \prod_{k\ne i}D(y_k)^{A_j}.
\end{equation}
The exclusion of $i$ from its coefficient is essential: the first
input law in Proposition~\ref{prop:tree-stability} has no inverse-margin
assumption.  Here is a proof.  For one level choose an index $r$ with
smallest reference margin $1-|y_r|$, and let $t$ be the smallest
margin among the other reference inputs.  If
$|x_k-y_k|\ge t/2$ for some $k\ne r$, the $k$th summand on the right,
already with exponent one, is at least $1/2$, since it contains
$D(y_r)\ge1/t$.  Increasing the constant bounds the left side by its
trivial upper bound $2$.

Otherwise every input other than $r$, along the straight segments
from $y$ to $x$, has margin at least $t/2$.  We claim that all partial
derivatives of $F_1$ there are bounded by a fixed power of $t^{-1}$,
even when the $r$th input is anywhere in the closed disk.  Holding
the siblings fixed, put
\[
 S_0=V-e-i\eta-\sum_{k\ne r}T_e(x_k),\qquad
 p=S_0-1/\bar q.
\]
The pole in the unrestricted Green variable is $p$, with
$\Im p\le-\Im(1/\bar q)\le-c$ and $|p|\le C/t$.
After substituting $T_e(x_r)$ and clearing $1-x_r$, the denominator
is, up to a nonzero uniformly bounded factor,
$(q-p)-x_r(\bar q-p)$.  Its modulus is at least
\[
 |q-p|-|\bar q-p|
 =\frac{-4\Im q\Im p}{|q-p|+|\bar q-p|}\ge ct.
\]
Numerator coefficients are polynomially bounded in $t^{-1}$;
sibling derivatives have the same property because
$|T_e'(x_k)|\le Ct^{-2}$.  This proves the claim, including the
removable value $x_r=1$.  Integration along the segments proves
\eqref{eq:tree-sibling-bound} for $j=1$: for $i=r$ its product
contains an input with margin $t$, and for $i\ne r$ it contains $r$.
For the induction, apply the one-level estimate to the outputs of
the child subtrees and then apply it inside the subtree containing
the varied leaf.  By \eqref{eq:tree-margin-bound}, each inverse
margin arising in another subtree is bounded by a sum of its leaf
inverse margins.  Since $D\ge1$ and the number of leaves is fixed,
these sums and their powers are bounded by a constant times a product
of sufficiently high powers.  None contains the varied leaf.  This
proves \eqref{eq:tree-sibling-bound} for every fixed $j$.
The one-level proof is unchanged for $d$ inputs at a full root,
and hence also for a fixed-depth recursion with that root degree.

\subsection{Mean and fluctuation contraction}

We next convert these deterministic estimates into a coupling bound.
Let $(X,Z)$ be a coupled pair of disk variables, let
$B=\|X-Z\|_2$ and $M=\EE(X-Z)$, and suppose both marginal $L^2$
norms are small.  Only $Z$ is assumed to satisfy an inverse-margin
bound $\EE D(Z)^p\le K$ for a sufficiently large fixed $p$.
Take independent copies $(X_i,Z_i)$ at the leaves, independent of
the internal potentials, and use identical internal potentials in
the two evaluations.  We claim that, for any prescribed
$\varepsilon>0$, sufficiently small marginal norms and sufficiently
small $h,w_0$ imply
\begin{align}
 \left|\EE\{F_1(X_\cdot)-F_1(Z_\cdot)\}-aM\right|
     &\le\varepsilon B,\label{eq:tree-mean-contraction}\\
 \|F_\ell(X_\cdot)-F_\ell(Z_\cdot)\|_2
     &\le\frac1{10}B+C_\ell|M|+\varepsilon B.
       \label{eq:tree-fluctuation-contraction}
\end{align}
The smallness threshold may depend on $K,p,\varepsilon,b,J,\ell$,
but not on the coupling.

We supply the decomposition behind this claim.  Fix $j=1$ or $\ell$,
write $N=b^j$, and, conditional on the internal potentials, set
$f_i(x)=F_j(0,\ldots,x,\ldots,0)$.  Then
\begin{equation}\label{eq:tree-single-decomposition}
 F_j(X_\cdot)-F_j(Z_\cdot)
 =\sum_{i=1}^{N}\bigl(f_i(X_i)-f_i(Z_i)\bigr)+R,
 \qquad \|R\|_2\le\varepsilon B,
\end{equation}
with any prescribed error after adjusting the thresholds.
Choose a small number $r>0$ and call pair $i$ bad when
$\max(|X_i|,|Z_i|)>r$.  If no pair is bad, telescoping over
coordinates and comparing each derivative with the corresponding
one-variable derivative gives $|R|\le\omega(r)\sum_i|X_i-Z_i|$,
where $\omega(r)\to0$.  If only pair $k$ is bad, the smoothness
on the one-unrestricted-input neighborhoods proved above gives
\[
 |R|\le\omega(r)|X_k-Z_k|+C\sum_{i\ne k}|X_i-Z_i|.
\]
For the latter terms we retain the indicator that pair $k$ is bad;
it is independent of $(X_i,Z_i)$.  Their $L^2$ norms are therefore
at most $CB\sqrt{\PP(k\text{ is bad})}$.

On the event of at least two bad pairs, use
\eqref{eq:tree-sibling-bound}, together with the bounded derivatives
of the $f_i$.  Each summand has the form
\[
 C|X_i-Z_i|\prod_{k\ne i}D(Z_k)^{A_j}
   \mathbf1_{\{\text{some }k\ne i\text{ is bad}\}}.
\]
Its $L^2$ norm is $B$ times the $L^2$ norm of the remaining factor.
For example, if $p\ge4A_j$, H\"older's inequality and independence
bound that norm by
\[
 C_{N,K}\PP(\text{some }k\ne i\text{ is bad})^{1/4}.
\]
The bad-pair probability tends to zero with
$\|X\|_2+\|Z\|_2$, by Markov's inequality.  First choose $r$ so
that $\omega(r)$ is small; then choose the marginal norms small.
Summing the finitely many estimates proves
\eqref{eq:tree-single-decomposition}.  This also explains why no
inverse-margin bound for $X$ is required.

The one-variable means satisfy, uniformly conditional on internal
potentials,
\begin{equation}\label{eq:tree-single-mean}
 \left|\EE\bigl[f_i(X_i)-f_i(Z_i)\mid\text{internal potentials}\bigr]
                  -f_i'(0)M\right|\le\varepsilon B.
\end{equation}
On $|X_i|,|Z_i|\le r$ this follows from derivative continuity.
Off that event, the global bounded derivative of $f_i$ and
Cauchy--Schwarz give an error at most
$CB\sqrt{\PP(i\text{ is bad})}$.  At one level the clean sum of
derivatives is $b(a/b)=a$; its perturbed value differs by
$O(h+w_0)$.  Equations~\eqref{eq:tree-single-decomposition} and
\eqref{eq:tree-single-mean} prove
\eqref{eq:tree-mean-contraction}.
At level $\ell$, the centered single-input differences are independent
conditional on the internal potentials.  By
\eqref{eq:tree-one-leaf-perturbed}, their sum has conditional centered
$L^2$ norm at most
$b^{\ell/2}2Cb^{-\ell}B<B/10$.  The sum of their conditional means
is bounded by $C_\ell|M|+\varepsilon B$, after making the individual
errors smaller by the fixed factor $N$.  Together with
\eqref{eq:tree-single-decomposition}, this proves
\eqref{eq:tree-fluctuation-contraction}.

\subsection{Smallness, inverse margins, and boundary continuity}

We can now prove Proposition~\ref{prop:tree-stability}.  First compare
the actual tree input with $Z=0$.  This comparison has inverse-margin
bound one and thus uses no unproved moment estimate for the tree.
As long as $B=\|Y\|_2$ is below the smallness threshold,
stationarity of the recursion and $\|F_j(0)\|_2\le C_j(h+w_0)$ give
\[
 |(1-a)\EE Y|\le\varepsilon B+C(h+w_0),\qquad
 B\le\tfrac1{10}B+C_\ell|\EE Y|+\varepsilon B+C_\ell(h+w_0).
\]
Since $|1-a|\ge c>0$, choose $\varepsilon$ small enough to absorb
the terms proportional to $B$.  It follows that
\begin{equation}\label{eq:tree-small-bootstrap}
 \|Y(e,\eta)\|_2\le C(h+w_0)
\end{equation}
whenever the initial smallness assumption holds.

Here is a starting point and continuation argument for that assumption.
At height $h$, the clean cavity function, obtained from its quadratic
equation, differs from $q(e)$ by $O(h)$ uniformly on $J$.  The resolvent
identity for adding the bounded potential changes a cavity entry by
at most $w_0/h^2$.  The map $X_e$ has uniformly bounded derivative
on $\HH$, so the starting disk norm is at most $C(h+w_0/h^2)$.
Choose $h$ small and then $w_0/h^2$ small enough to start strictly
inside the threshold and to make the right side of
\eqref{eq:tree-small-bootstrap} less than half that threshold.
For each fixed $e,w$, the tree resolvent, and hence its disk-valued
$L^2$ norm, is continuous for $\eta>0$.  A first crossing of the
threshold on decreasing $\eta$ would contradict the improved bound.
This proves \eqref{eq:tree-small-bootstrap} throughout $(0,h]$,
uniformly in $e,w$.  Decreasing the parameters makes the resulting
norm as small as any prescribed fixed number.

Next fix a finite moment order $p$.  At every strictly positive height
all inverse margins have finite deterministic bounds: for the bounded
tree operator,
$|g|\le\eta^{-1}$ and
$\Im g\ge\eta/(C^2+\eta^2)$ on our bounded energy range.
It is therefore legitimate to absorb an as yet nonuniform moment.
Let $N=b^\ell$ and use the radius $r$ for which
\eqref{eq:tree-one-leaf-perturbed} has a uniformly interior image
when at most one leaf is outside $|x|\le r$.  On that event
$D(F_\ell(Y_\cdot))^p\le C_p$.  On its complement,
\eqref{eq:tree-margin-bound} iterated through $\ell$ levels gives
$D(F_\ell(Y_\cdot))^p\le C_p\sum_iD(Y_i)^p$.
For each summand the complement implies that some other leaf is
outside $|x|\le r$.  Independence and recursion stationarity yield
\[
 \EE D(Y)^p
 \le C_p+C_pN(N-1)\PP(|Y|>r)\,\EE D(Y)^p.
\]
By \eqref{eq:tree-small-bootstrap}, the probability coefficient can
be made less than $1/(2C_pN(N-1))$.  Absorption gives a uniform
bound $2C_p$.  Taking the largest of finitely many required orders
handles all of them, with constants that remain bounded as the
parameters decrease further.

We have now obtained smallness without moments and then obtained moments
from smallness.  Fix those moment constants and apply the coupling
estimates with $Z$ having the tree law.  Choose an optimal coupling
of $P$ and that law, so $B=W_2(P,\Law(Y))$ and
$M=\EE(X-Y)$.  The first approximate recursion and
\eqref{eq:tree-mean-contraction} imply
$|(1-a)M|\le u+\varepsilon B$, hence
$|M|\le C(u+\varepsilon B)$.  The second approximate recursion,
the coupling through identical internal potentials, and
\eqref{eq:tree-fluctuation-contraction} give
\[
 B\le u+\tfrac1{10}B+C_\ell|M|+\varepsilon B.
\]
Absorbing again proves \eqref{eq:tree-law-stability}.  The coupling
smallness threshold $\delta_*$ is now fixed.  Decreasing $h,w_0$
once more makes the actual tree size arbitrarily smaller than
$\delta_*$ without enlarging its moment bounds.  This is the
parameter order asserted in the proposition.

Finally compare the two actual tree variables on the same potentials
at parameters $(e,\eta)$ and $(e',\eta')$, both at positive height.
Use their own coordinates $X_e,X_{e'}$, and write the coupled variables
as $Y,Y'$.  Their leaf pairs are independent across branches, and
recursing preserves this particular coupling.  When the second
evaluation is temporarily made using the first parameters, the
coupling estimates apply with $B=\|Y-Y'\|_2$, rather than merely
with a transportation distance.  The parameter replacement error in
either $F_1$ or $F_\ell$ has $L^2$ norm at most a modulus
$r_0(|e-e'|+|\eta-\eta'|)$ tending to zero uniformly.  To see this,
restrict all finitely many leaves to $D(Y_i')\le R$; on that compact
subset the recursion maps are uniformly continuous in their
parameters, including at height zero.  The complementary probability
is at most $NK_pR^{-p}$, and the disk-valued difference is bounded
by $2$.  First choose $R$ large, and then make the parameter distance
small.  Thus the preceding mean and norm inequalities now have
$r_0$ in place of $u$, and imply
\[
 \|Y(e,\eta)-Y(e',\eta')\|_2
 \le C r_0(|e-e'|+|\eta-\eta'|).
\]
Completeness of $L^2$ gives the asserted continuous extension.
Fatou's lemma along almost surely convergent subsequences preserves
all inverse-margin bounds, so the limit lies in the open disk.
Each boundary variable is a limit in probability of measurable
functions of the potentials and is therefore measurable in that
environment.  Since the finitely many child limits lie in the open
disk, continuity of the recursion also passes
\eqref{eq:tree-recursion} to height zero.  These identities hold
simultaneously at all vertices for each fixed energy, by taking the
intersection of countably many probability-one events.
This finishes the proof of
Proposition~\ref{prop:tree-stability}.

\subsection{Spatial covariances and the density of states}

At the full root, let $g_1,\ldots,g_d$ denote the incoming cavity
functions, and put
\[
 m_o(\zeta)=\left(V_o-\zeta-\sum_{i=1}^d g_i(\zeta)\right)^{-1}.
\]
The next lemma makes explicit which spatial estimates are required
later.  A star consists of the root and its $d$ neighbors.  By a
retained star we mean the coordinate set obtained by omitting one
neighbor, while keeping the root and the other $d-1$ neighbors.

\begin{lemma}[Tree covariance and logarithmic integrability]
\label{lem:tree-log}
Fix a finite exponent $p_0\ge1$ before choosing the disorder threshold.
Taking the finite order $p_*$ in
Proposition~\ref{prop:tree-stability} sufficiently large, the full-root
Green function and every entry of the star Green matrix extend
continuously in $L^{p_0}$ to $J\times[0,h]$ on the common environment.
Their boundary values have strictly positive
root imaginary part.  If $K(\zeta)$ is the imaginary part of the
Green matrix compressed to an edge or a retained star, then
$K(e+i0)$ is positive definite almost surely.  Uniformly for
$e\in J$ and $0\le\eta\le h$, the variables
$\log\det K(e+i\eta)$ are uniformly integrable, and they converge
in $L^1$ when $(e,\eta)$ converges.  These assertions require only
a finite collection of inverse-margin moments of spatial cavity
variables, independent of the number of colors in a Gaussian field.
\end{lemma}

\begin{proof}
Put $M_0=2+\sum_{i=1}^dD(Y_i)$.  The comparison
\eqref{eq:tree-hyperbolic-comparison} gives polynomial bounds in
$M_0$ for $|g_i|$, $|g_i|^{-1}$, and $(\Im g_i)^{-1}$.
For $A=V_o-\zeta-\sum_i g_i$, we have $|A|\le CM_0$ and
$-\Im A\ge cM_0^{-1}$.  Hence $|m_o|$, $|m_o|^{-1}$, and
$(\Im m_o)^{-1}$ are also bounded by fixed powers of $M_0$.
In particular $\Im m_o(e+i0)>0$ almost surely.

The inverse star Green matrix and its inverse have the explicit forms
\begin{equation}\label{eq:tree-star-matrix}
 S=\begin{pmatrix}
 V_o-\zeta&1&\cdots&1\\
 1&1/g_1&&0\\
 \vdots&&\ddots&\\
 1&0&&1/g_d
 \end{pmatrix},\qquad
 S^{-1}=\begin{pmatrix}
 m_o&(-m_og_i)_i\\
 (-m_og_i)_i^{\mathsf T}&(\delta_{ij}g_i+g_im_og_j)_{ij}
 \end{pmatrix}.
\end{equation}
Thus $\|S\|+\|S^{-1}\|\le CM_0^C$.  Also
\[
 -\Im S=\diag\left(\eta,
       \frac{\Im g_1}{|g_1|^2},\ldots,
       \frac{\Im g_d}{|g_d|^2}\right),
 \qquad
 \Im S^{-1}=(S^{-1})^*(-\Im S)S^{-1}.
\]
Each neighbor diagonal in $-\Im S$ is at least $cM_0^{-C}$.
For a vector $x$ supported on a retained star, let the omitted
neighbor be $r$ and put $y=S^{-1}x$.  The unused row says
$0=y_o+y_r/g_r$, so $|y_o|\le CM_0^C|y_r|$.
Consequently
\[
 |x|=|Sy|\le CM_0^C|(y_1,\ldots,y_d)|,
 \qquad
 x^*(\Im S^{-1})x\ge cM_0^{-C}|x|^2.
\]
Together with the upper bound for $S^{-1}$, this bounds every
eigenvalue of the retained covariance between $cM_0^{-C}$ and
$CM_0^C$.  The same bounds hold for an edge covariance: choose a
retained star that contains the edge and restrict its quadratic
form.  This is possible since $d\ge3$.  The bounds hold equally
at height zero by the displayed formulas and the almost sure
positivity of all incoming imaginary parts.

All formulas are continuous when their inputs lie in $\DD$.
Proposition~\ref{prop:tree-stability} therefore gives convergence in
probability.  The polynomial bounds and sufficiently high, but
finitely many, inverse-margin moments give the claimed $L^{p_0}$
convergence by uniform integrability.  Moreover
\[
 |\log\det K|\le C(1+\log M_0).
\]
Any fixed positive moment of $M_0$ bounds, for example, the second
moment of this logarithm.  The logarithms are consequently uniformly
integrable and converge in $L^1$.  Tensoring $K$ with the identity
on $m$ colors multiplies $\log\det K$ by $m$; it does not require
any larger spatial moment order.
\end{proof}

\begin{corollary}[Continuous positive density of states]
\label{cor:tree-density}
For the above smallness choices, the restriction of $\nu_{d,w}$ to
the interior of $J$ has the continuous strictly positive density
\begin{equation}\label{eq:tree-density}
 \rho_{d,w}(e)=\frac1\pi\EE\Im m_o(e+i0).
\end{equation}
In particular this holds on an open neighborhood of
$I_{d,\kappa}$ when that interval lies in the interior of $J$.
\end{corollary}

\begin{proof}
By the definition of the density-of-states measure,
$\EE m_o(\zeta)=\int(t-\zeta)^{-1}\,d\nu_{d,w}(t)$.
Lemma~\ref{lem:tree-log} gives $L^1$ continuity of $m_o$ on
$J\times[0,h]$, uniformly on this compact parameter set.  Thus
$\pi^{-1}\EE\Im m_o(e+i\eta)$ converges uniformly on $J$ to
the continuous function in \eqref{eq:tree-density}.  For every
continuous compactly supported $f$ in the interior of $J$, Poisson
kernel approximation and Fubini's theorem give
\[
 \int f\,d\nu_{d,w}
 =\lim_{\eta\downarrow0}\frac1\pi
       \int f(e)\EE\Im m_o(e+i\eta)\,de
 =\int f(e)\rho_{d,w}(e)\,de.
\]
This identifies the entire restricted measure, including the absence
of a singular component.  Since the root imaginary part is strictly
positive almost surely, its finite expectation is strictly positive
at every $e\in J$.
\end{proof}
